\documentclass[
10pt, % Main document font size
a4paper, % Paper type, use 'letterpaper' for US Letter paper
oneside, % One page layout (no page indentation)
%twoside, % Two page layout (page indentation for binding and different headers)
headinclude,footinclude, % Extra spacing for the header and footer
BCOR5mm, % Binding correction
]{scrartcl}
\input{stru/structure.tex}
\hyphenation{Fortran hy-phen-ation} 

\title{\normalfont{FTML practical session 14}} 
\author{\spacedlowsmallcaps{}
} 
\date{\today}
\begin{document}
\renewcommand{\sectionmark}[1]{\markright{\spacedlowsmallcaps{#1}}}
\lehead{\mbox{\llap{\small\thepage\kern1em\color{halfgray} \vline}\color{halfgray}\hspace{0.5em}\rightmark\hfil}} 
\pagestyle{scrheadings}
\maketitle 
\setcounter{tocdepth}{3} 

\tableofcontents

\section{\large\color{Blue}Lower bound on the performance of gradient-based algorithms}

In previous sessions, we have discussed several upper bounds on the performance
of machine learning methods. 

For instance, we have seen that for a strongly
convex smooth function $f$ defined on $ \mathbb{R}^d$ and with real values,
the convergence of gradient descent to the minimizer $x^*$ of
$f$ can be upper bounded. More formally, we have seen that if
$(x^t)_{t\in N\mathbb{n} }$ is the sequence of iterates, then

\begin{equation}
    ||x^t-x^*|| = \mathcal{O}(\alpha^t) 
\end{equation}

with $0\leq \alpha<1$. If $f$ is only convex, and not necessary strongly
convex, then we have a weaker upper bound:

\begin{equation}
    f(x^t)-f(x^*) = \mathcal{O} ( 1/t)
\end{equation}

Having an upper bound guarantees that the algorithm converges \textbf{{faster}}
than the bound.

In this session we present the problem of obtaining \textbf{{lower}} bounds on the
performance of machine learning methods. This means proving that an algorithm
can not be guaranteed to converge too fast to the minimizer of $f$. We will
formally define what this means in section \ref{sec:minimax}, as it is not as straightforward as for upper
bounds.

There are two kinds of lower bounds: statistical lower bounds and optimization lower bounds, which we will focus on.


Let $ \mathcal{F} $ be the space of functions $f: \mathbb{R}^d \rightarrow \mathbb{R} $ that
\begin{itemize}
    \item are convex and differentiable
    \item and have a minimizer $x_f^*$
    \item are $L$-smooth, with $L\in \mathbb{R}_+$ (the gradient is $L$-Lipshitz continuous).
\end{itemize}

\subsection{\large\color{MidnightBlue}First-order methods}

We consider optimization algorithms that are based on linear combinations of local estimations of the gradient of $f$. These methods are called \textbf{{first-order methods.}} GD is an example of first-order method. Formally, this means that each iterate $x^t$ verifies

\begin{equation}
    \label{eq:first}
    x^t\in x^0+ \text{span}\{\nabla f(x^0), \dots, \nabla f(x^{t-1}\}
\end{equation}

"span" means "espace engendré par".

\subsection{\large\color{MidnightBlue}Nesterov acceleration (AGD)}

Nesterov acceleration (also called accelerated gradient descent (AGD)), is another first-order method, slightly different than GD. It is possible to prove that with Nesterov acceleration, the rate of convergence is $ \mathcal{O}(1/t^2)$, instead of $ \mathcal{O}(1/t)$.

\url{https://blogs.princeton.edu/imabandit/2013/04/01/acceleratedgradientdescent/} 

\subsection{\large\color{MidnightBlue}Minimax lower bound}
\label{sec:minimax}

It is possible to show that AGD is optimal among first order methods. But first, we have to define what this means. This optimality is defined as a minmax problem. Indeed, we are interested in finding the algorithm that has the best worst-case performance. It is necessary to restrict algorithms to a relevant class of algorithms $ \mathcal{A} $ (here, first order methods), and the functions to a relevant class of functions $ \mathcal{F} $ (for instance the set of convex, $L$-smooth functions on $\mathbb{R}^d$). 

For $t\in \mathbb{N}$, we look for an algorithm $a\in \mathcal{A}$ defined by

\begin{equation}
    a = \argmin_{a\in \mathcal{A} } \max_{f\in \mathcal{F} } \Big[||x^t_a-x_f^*||\Big]
\end{equation}

where $x^t_a$ is the iterate returned by algorithm $a$ at iteration $t$.

\subsection{\large\color{MidnightBlue}Optimality of Nesterov acceleration}

If we manage to show that for any algorithm $a\in \mathcal{A}$, there exists a function $f\in \mathcal{F} $, such that 

\begin{equation}
    f(x^t)-f(x^*)\geq \frac{C}{t^2} ||x^0-x_f^*||^2
\end{equation}

where $C$ is independent on the function, then this will prove that Nesterov acceleration is optimal in $ \mathcal{A} $, as this means that for any algorithm in $ \mathcal{A} $, there exists a function, for which the iterates computed by the algorithm $a$ converge slower than $ \frac{C}{t^2} $ to its minimizer.

\subsection{\large\color{MidnightBlue}Hard function}

Without loss of generality, we assume that for all algorithms in $ \mathcal{A}$, the initialization is $x^0=0\in \mathbb{R}^d$. 

We consider $t\leq\frac{d-1}{2} $ and $f$ defined by

\begin{equation}
    f(x) = \frac{L}{4} \Big( \frac{1}{2} x_1^2 + \frac{1}{2} \sum^{2t}_{i=1} (x_i-x_{i+1})^2+ \frac{1}{2} x_{2t+1}^2-x_1 \Big)
\end{equation}

\exercice{\textbf{}}Show that $f$ is convex.
\\

We admit that $f$ is $L$-smooth, and that the minimizer of $f$, $x^*$, is

$$
x_i^*  : \left\{
    \begin{array}{ll}
	1  - \frac{i}{2t+2} \text{ for } 1\leq i \leq 2t+1 & \\
	0 \text{ for } i\geq 2t+2
    \end{array}
\right.
$$

and that the minimum value is

\begin{equation}
    f^* = \frac{L}{8} \big( \frac{1}{2t+2} -1\big)
\end{equation}

We consider another utility function $g$ defined as

\begin{equation}
    g(x) = \frac{L}{4} \Big( \frac{1}{2} x_1^2 + \frac{1}{2} \sum^{t-1}_{i=1} (x_i-x_{i+1})^2+ \frac{1}{2} x_{t}^2-x_1 \Big)
\end{equation}

We also admit that the minimum value of $g$, noted $g^*$, is

\begin{equation}
    g^* = \frac{L}{8} \big( \frac{1}{t+1} -1\big) 
\end{equation}

\exercice{\textbf{}} Compute the gradient of $f$ for any point $x\in \mathbb{R}^d$.
\\


\exercice{\textbf{}} Show that for all $l\leq d$, the components $j$ with $j\geq l+1$ are null for the iterate $x^l$, for any first-order method that produces iterates with gradients of $f$ (as defined in \ref{eq:first}). In particular, this is true for $l=t$ and for the iterate $t$, we have

\begin{equation}
    x^t_{t+1} = x^t_{t+2} = \dots = x^t_d = 0   
\end{equation}

\exercice{\textbf{}}Show that $f(x^t)=g(x^t)$.
\\

Hence, $f(x^t)\geq g^*$ and 

\begin{equation}
    \frac{f(x^t)-f^*}{||x^0-x^*||^2} \geq \frac{g^*-f^*}{||x^0-x^*||^2} = \frac{g^*-f^*}{||x^*||^2}
\end{equation}

\exercice{\textbf{}}Show that

\begin{equation}
    ||x^*||^2\leq \frac{2t+2}{3} 
\end{equation}


\exercice{\textbf{}}Conclude.

\subsection{\large\color{MidnightBlue}Discussion}

By essence, considering the worst-case performance (the hardest function to optimize) is pessimistic. Some functions optimized wit a first-order method by converge faster than $ \mathcal{O}(1/t^2)$. We have seen this behavior for a least-squares setting in a previous session. 


\end{document}
